\section{Conclusion}

In this paper, we propose an IQN-based local path planner for sensor-based Unmanned Surface Vehicle navigation in marine environments with no prior knowledge of the current flow field and obstacles, and further enhance safety via a CVaR-style risk distortion with a distance-adaptive risk mechanism. Compared with a DQN-based planner, our method shows more stable learning behavior.
Compared with state-of-the-art DRL baselines, our method dominates the performance trade-off in dense, dynamic environments. It attains a perfect \textbf{100\% success rate}, closing the safety gap seen in the highly efficient PPO (98\%) without sacrificing its speed or energy advantages ($\sim$21.50\,s, 14.95\,J). While classical planners (APF) provide a strong safety baseline (95\%), they do so at the cost of severe inefficiency in complex flows. The proposed distance-adaptive risk mechanism successfully bridges this divide, providing a robust, real-time solution for mission-critical marine autonomy.
